% TOP_PROBLEM: {"id": 1281, "title": "Four Exponentials Conjecture", "category_id": 1, "category": {"id": 1, "name": "number_theory", "display_name": "Number Theory", "description": "Properties of integers, prime numbers, Diophantine equations.", "slug": "number-theory", "order_index": 1, "created_at": "2026-07-31T15:26:25.670Z"}, "status": "open"}
% FIRST_POSTED: 2026-09-27
% !TeX program = lualatex
% ATTEMPT_STATUS: unresolved
\documentclass[11pt]{article}
\usepackage[margin=25mm]{geometry}
\usepackage{fontspec}
\setmainfont{Latin Modern Roman}
\usepackage{amsmath,amssymb,mathtools}
\usepackage{unicode-math}
\setmathfont{Latin Modern Math}
\usepackage{xurl,hyperref}
\hypersetup{colorlinks=true,urlcolor=blue}
\setcounter{secnumdepth}{0}
\setlength{\parindent}{0pt}
\setlength{\parskip}{5pt}
\setlength{\emergencystretch}{3em}
\begin{document}
\section*{\texorpdfstring{Four Exponentials Conjecture}{Four Exponentials Conjecture}}\label{1281}
\subsection{Definitions and mathematical statement}
Let $x_1,x_2,y_1,y_2\in{\mathbb{C}}$ with $x_1,x_2$ linearly independent over ${\mathbb{Q}}$ and $y_1,y_2$ linearly independent over ${\mathbb{Q}}$. Thus each rational relation within either pair is trivial. The four exponentials conjecture asserts
\[
 \{e^{x_1y_1},e^{x_1y_2},e^{x_2y_1},e^{x_2y_2}\}
 \not\subset\overline{{\mathbb{Q}}}.
\]
Here $\overline{{\mathbb{Q}}}$ denotes the algebraic complex numbers, so at least one of the four values must be transcendental. The pairs need not consist of real numbers, and no independence between the two pairs is assumed. The assertion does not identify which value is transcendental or demand algebraic independence of the four values.
\par\smallskip\textit{Formulation and concept references:} \href{https://mathworld.wolfram.com/FourExponentialsConjecture.html}{[S1]}.

\subsection{Short English statement}
For two pairs of rationally independent complex numbers, must at least one exponential of a cross-product be transcendental?
\subsection{Research attempt}
\textbf{Status: UNRESOLVED.} No unconditional proof of the four exponentials conjecture is obtained. This attempt proves some algebraic-ratio cases using classical transcendence theorems, gives a detailed conditional implication from Schanuel's conjecture, and isolates the rank-one matrix obstruction which those arguments leave unresolved. Complex logarithm branches are kept explicit throughout.

\paragraph{Source audit.}
The original formulation is retained above. For the distinction between four and six exponentials, and for the implication from algebraic independence of logarithms, I use Waldschmidt's author-hosted survey [E1], consulted on 27 September 2026. It explicitly distinguishes the conjectural four-exponentials statement from the proved six-exponentials theorem. The linear-space argument below expands the elementary rank-one geometry behind that conditional implication; it is not a new unconditional transcendence theorem.

\paragraph{1. Normalize without strengthening the assumptions.}
Independence implies that all $x_i,y_j$ are nonzero. Set
\[
 u=x_1y_1,\qquad r=x_2/x_1,\qquad s=y_2/y_1.
\]
The exact hypotheses become $u\ne0$, $r\notin\mathbb Q$, and $s\notin\mathbb Q$. The four numbers are
\[
 e^u,\quad e^{su},\quad e^{ru},\quad e^{rsu}.
\]
Neither $r$ nor $s$ is required to be real; there is no hypothesis of independence between all four original parameters. Scaling the $x$'s by a nonzero complex number and the $y$'s by its inverse does not change the four products. Thus introducing $u,r,s$ loses no information relevant to the assertion.

The exponent matrix
\[
 M=u\begin{pmatrix}1&s\\r&rs\end{pmatrix}
\]
has complex rank one, while neither its two rows nor its two columns are rationally proportional. The conjecture says that a matrix of this kind cannot have every entry among the logarithms of nonzero algebraic numbers. This is a condition on exact logarithms, not simply on the algebraic values of four exponentials taken without their branches.

Both independence hypotheses matter. Take $x_1=1,x_2=2,y_1=\log2,y_2=\log3$, with real logarithms. The two $y$'s are rationally independent: an integral relation would give $2^a3^b=1$, and unique factorization gives $a=b=0$. Nevertheless the four exponentials are $2,3,4,9$. The failure is precisely the rational dependence of the $x$'s. Transposing this example tests the other independence assumption.

\paragraph{2. Two genuinely unconditional subcases.}
The Hermite--Lindemann theorem says that the exponential of a nonzero algebraic number is transcendental. Therefore the desired conclusion holds if even one of $u,su,ru,rsu$ is algebraic. All four products are nonzero, so no exception involving $e^0$ arises.

There is another useful subcase. Suppose $r$ is algebraic and irrational. If $e^u$ is transcendental, there is nothing to show. Otherwise write $\alpha=e^u\in\overline{\mathbb Q}^{\times}$. If $\alpha\ne1$, the Gelfond--Schneider theorem applied with the logarithm $u$ of $\alpha$ says that
\[
 e^{ru}=\alpha^r
\]
is transcendental. Here the theorem is used in its form covering every value of the algebraic power; no principal-branch assumption is made.

The case $\alpha=1$ deserves separate treatment because it is excluded from that statement of Gelfond--Schneider. Then $u=2\pi i k$ for a nonzero integer $k$. Using the logarithm $\pi i$ of $-1$, write
\[
 e^{ru}=e^{(2kr)\pi i}=(-1)^{2kr}.
\]
The exponent $2kr$ is again algebraic and irrational, and the base $-1$ is algebraic and different from zero and one. Gelfond--Schneider applies. The same reasoning with columns proves the conclusion when $s$ is algebraic and irrational.

Consequently a counterexample would have both $r$ and $s$ transcendental, and every exponent entry transcendental. This necessary condition is far from sufficient for either constructing or excluding a counterexample. In particular, applying Gelfond--Schneider when the exponent itself is transcendental is not justified; algebraic bases can certainly have algebraic values at some transcendental exponents.

\paragraph{3. A rational rank-one linear-space lemma.}
We need an elementary algebraic fact. Let $W$ be a rational vector subspace of $2\times2$ matrices such that every matrix in $W$ has determinant zero. If $W\ne0$, then either all its columns lie in one fixed rational line or all its rows lie in one fixed rational line.

Choose a nonzero member $A=ab^{\mathsf T}$, where $a,b\in\mathbb Q^2$ are nonzero; this factorization exists because $A$ has rational entries and rank one. Any other nonzero $B=cd^{\mathsf T}$ has
\[
 \det(A+B)=\det(a,c)\det(b,d).
\]
The sign convention in this formula is immaterial for vanishing, and direct expansion verifies the displayed convention. Thus either $c$ is proportional to $a$ or $d$ is proportional to $b$.

If every $B$ has column line $\mathbb Qa$, the first alternative holds globally. Otherwise choose $B$ whose column line differs from that of $A$. Its row line must be $\mathbb Qb$. If some $C$ had row line different from $\mathbb Qb$, the preceding dichotomy would force its column line to be $\mathbb Qa$. Then $B+C$ would have both independent column factors and independent row factors, so its determinant would be nonzero. That contradicts $B+C\in W$. Hence all rows lie in $\mathbb Qb$. Zero members do not affect either conclusion, proving the lemma.

Equivalently, a rational linear subspace contained in the determinantal quadric has a common rational left kernel or a common rational right kernel. The rational qualification is essential for the next step: a common arbitrary complex line would merely restate that the original matrix has rank one and would not contradict rational independence.

\paragraph{4. A complete conditional deduction from Schanuel.}
Assume Schanuel's conjecture in the following usual form: for rationally independent complex numbers $z_1,\ldots,z_d$,
\[
 \operatorname{trdeg}_{\mathbb Q}
 \mathbb Q(z_1,\ldots,z_d,e^{z_1},\ldots,e^{z_d})\ge d.
\]
Suppose toward a contradiction that all four exponentials in the problem are algebraic. Choose a rational basis $z_1,\ldots,z_d$ for the span of the four entries of $M$, choosing the basis from those entries. Then every $e^{z_j}$ is algebraic. Schanuel therefore says that the $d$ numbers $z_j$ themselves are algebraically independent.

Write $M=\sum_jz_jA_j$ with rational matrices $A_j$. Its determinant is zero, so
\[
 P(z_1,\ldots,z_d)=0,\qquad
 P(T_1,\ldots,T_d)=\det\Bigl(\sum_jT_jA_j\Bigr).
\]
Algebraic independence forces $P$ to be the zero polynomial. Hence every member of the rational span $W$ of the $A_j$ is singular. This span is nonzero because $M\ne0$. Apply the lemma.

If $W$ has a common rational column line spanned by $(a_1,a_2)^{\mathsf T}$, the same is true of its complex linear combinations, in particular $M$. Its first column is $(u,ru)^{\mathsf T}$ with $u\ne0$, so $a_1\ne0$ and $r=a_2/a_1\in\mathbb Q$. If $W$ has a common rational row line, the first row $(u,su)$ similarly gives $s\in\mathbb Q$. Either conclusion contradicts the hypotheses. This proves the four-exponentials statement under Schanuel.

The implication needs only the special consequence that rationally independent logarithms of algebraic numbers are algebraically independent. This weaker conjectural input is still unproved in the required generality. Omitting the sentence invoking it would leave an invalid step: rational independence alone never licenses concluding that a multivariate determinant polynomial vanishes identically from its vanishing at one tuple.

\paragraph{5. Why six exponentials cannot simply be restricted.}
The six-exponentials theorem supplies a transcendental member of a $2\times3$ array when the two row parameters and the three column parameters are separately rationally independent. Starting with the present $y_1,y_2$, choosing $y_3=y_1+y_2$ keeps the new exponentials algebraic if the old ones are algebraic, but destroys the required independence of the three column parameters. Choosing a genuinely independent $y_3$ fixes that hypothesis, but then the theorem may locate its only transcendental value in the new third column.

This is a precise quantifier obstruction. A theorem asserting that at least one of six entries is transcendental does not assert that an entry in any prescribed four-entry subarray is transcendental. Multiplicative identities among the exponential values do not remedy the missing independence: they often arise exactly from the rational relations which prevent application of the six-exponentials theorem.

There is also a branch obstruction to a proposed counterexample. Choosing four algebraic numbers $\alpha_{ij}$ gives logarithms of the form
\[
 L_{ij}+2\pi i k_{ij},\qquad k_{ij}\in\mathbb Z,
\]
for any initially chosen logarithms $L_{ij}$. One must make their determinant exactly zero while preserving the rational independence of the two rows and two columns. A numerical near-zero determinant cannot establish that equality, and changing one branch generally changes the determinant. The algebraic nature of the four $\alpha_{ij}$ imposes no automatic rank-one identity on their logarithms.

\paragraph{6. A small-looking consequence and its limitation.}
Take $x_1=\log2,x_2=\log3,y_1=1,y_2=t$, where $t$ is real and irrational. The first pair is rationally independent by unique factorization and the second by irrationality. Two entries are already the algebraic numbers 2 and 3. The conjecture therefore implies that at least one of $2^t,3^t$ is transcendental. If $t$ is algebraic irrational, Gelfond--Schneider already proves this, separately for each base. The remaining case permits transcendental $t$ and is not reached by that theorem.

The offline diagnostics are limited to exact algebra. They enumerate small rational rank-one matrices and test the determinant identity and common-line conclusion for singular linear spans. They also check the algebraic substitutions in the normalization. There is no numerical ``test for transcendence'', no certification of rational independence from a floating-point logarithm, and no inference about all complex parameters from a finite search.

The missing unconditional ingredient can now be stated sharply: show that a rank-one matrix of logarithms of algebraic numbers has a rational dependence between its rows or between its columns, without assuming the algebraic independence consequence of Schanuel. The special algebraic-ratio cases remove some configurations, while the conditional argument explains the intended obstruction. Neither provides that last transcendence statement.

\subsection{Sources}
\begin{itemize}
\item[S1]Four Exponentials Conjecture (MathWorld).
\url{https://mathworld.wolfram.com/FourExponentialsConjecture.html}.
\textit{Formulation source.}
\item[E1]Michel Waldschmidt, The Four Exponentials Problem and the Schanuel Conjecture, author manuscript dated 26 May 2022, especially Sections 3--5.
\url{https://webusers.imj-prg.fr/~michel.waldschmidt/articles/pdf/FourExponentialsSchanuel.pdf}.
\textit{Added primary source: Precise four- and six-exponentials statements and the conditional rank argument. Consulted 27 September 2026.}
\end{itemize}
\end{document}
